\documentclass[11pt]{article}

\usepackage[margin=1in]{geometry}
\usepackage{amsmath}
\usepackage{amssymb}
\usepackage{booktabs}
\usepackage{hyperref}
\usepackage{listings}
\usepackage{xcolor}

\hypersetup{
  colorlinks=true,
  linkcolor=blue!55!black,
  citecolor=blue!55!black,
  urlcolor=blue!55!black
}

\lstset{
  basicstyle=\ttfamily\small,
  breaklines=true,
  frame=single,
  columns=fullflexible
}

\title{A Browser-Based Wide-Field Star Calibration Tool for AIDA Camera Models}
\author{Juha Vierinen \and Bjorn Gustavsson \and Codex}
\date{Draft for AIDA Browser Calibrator v0.2.0}

\begin{document}
\maketitle

\begin{abstract}
We describe a browser-based calibration tool for fitting wide-field camera
models to images of the night sky. The application combines manual star
pairing, automatic bright-star detection, blind asterism matching, robust
nonlinear optimization, residual inspection, and code export in a self-contained
HTML and JavaScript interface. The current implementation is designed for
wide-field and all-sky cameras with fields of view above roughly 30 degrees and
stars brighter than visual magnitude 7. It supports the parametric AIDA optical
models used in MATLAB tooling and a Brown-Conrady model for ordinary phone
cameras. The tool emphasizes practical calibration under imperfect observing
conditions, including aurora, buildings, image noise, nonlinear lens effects,
and spurious star detections.
\end{abstract}

\section{Introduction}

Wide-field optical calibration is needed when image pixels must be mapped to
sky directions. Examples include all-sky auroral imaging, meteor cameras,
planetarium-style camera alignment, and phone-camera photographs of bright
stars. Traditional calibration workflows often combine several separate tools:
image viewers, star catalogs, custom fitting scripts, and hand-written export
code. The AIDA browser calibrator packages these steps into one application.

The application is intentionally local and transparent. Images are loaded in the
browser. The bright-star catalog is embedded in JavaScript. The lens model,
projection, star detection, asterism matching, fitting, residual analysis, and
export routines are all implemented client-side. When served by the local
development server, the application can also save calibrated test cases for
regression testing.

\section{User Workflow}

The workflow is centered on matching image detections to catalog stars. The user
loads an image, verifies the UTC timestamp and observer location, selects an
optical model, and roughly aligns the sky projection. The application supports
manual pairing through a keyboard-modified click workflow. Holding \texttt{S}
and clicking near an image star opens a local density estimate; the centroid is
refined at subpixel resolution and can then be paired with a catalog star.

The ``I'm feeling lucky'' path automates the same idea. It detects bright image
features, attempts a blind spherical asterism match, seeds the selected optical
model, expands the number of candidate matches, fits the lens parameters, and
rejects obvious outliers. The user can paint bad star-finder regions with
\texttt{M}+click. These marked regions remove yellow automatic detections from
future automatic matching and are saved with submitted test-case metadata.

\section{Coordinate Systems}

The application uses true zero-based image pixel coordinates in the browser. A
catalog star is first converted from right ascension and declination to local
azimuth and zenith angle using the image timestamp and observer position. Let
\(\alpha_\star\) be right ascension, \(\delta_\star\) declination, \(\phi\) the
observer latitude, and \(H\) the local hour angle. The altitude is
\begin{equation}
  h = \sin^{-1}\left(\cos H \cos\delta_\star \cos\phi
      + \sin\delta_\star \sin\phi\right),
\end{equation}
and the zenith angle is
\begin{equation}
  z = \frac{\pi}{2} - h.
\end{equation}
The local sky unit vector used by the AIDA camera model is
\begin{equation}
  \mathbf{e} =
  \begin{bmatrix}
    \sin z \sin A \\
    \sin z \cos A \\
    \cos z
  \end{bmatrix},
\end{equation}
where \(A\) is azimuth.

\section{Camera Rotation and Projection}

The browser implementation follows the AIDA convention of a three-angle camera
rotation. The rotation matrix is
\begin{equation}
  R = R_y(\alpha) R_x(\beta) R_z(\gamma).
\end{equation}
The local sky vector is transformed into camera coordinates by
\begin{equation}
  \mathbf{s} = \mathbf{e} R,
  \qquad
  \mathbf{s} = (s_1, s_2, s_3).
\end{equation}
The radial angle from the optical axis is
\begin{equation}
  \theta = \tan^{-1}\left(\frac{\sqrt{s_1^2+s_2^2}}{s_3}\right).
\end{equation}

For AIDA \texttt{optmod 2}, the radial mapping is
\begin{equation}
  r = \sin(a_r \theta),
\end{equation}
and the normalized pixel coordinates are
\begin{align}
  u &= f_1 \frac{s_1}{\sqrt{s_1^2+s_2^2}} r + \frac{1}{2} + d_u, \\
  v &= f_2 \frac{s_2}{\sqrt{s_1^2+s_2^2}} r + \frac{1}{2} + d_v.
\end{align}
The browser pixel coordinates are then
\begin{equation}
  x = u W - 1, \qquad y = v H - 1,
\end{equation}
where \(W\) and \(H\) are image width and height. The \(-1\) term converts from
the historical MATLAB one-based convention to browser zero-based coordinates.

\section{Supported Optical Models}

The current browser supports AIDA parametric optical models \texttt{1},
\texttt{2}, \texttt{3}, \texttt{4}, \texttt{5}, and \texttt{12}. These share the
parameter vector
\begin{equation}
  \mathbf{p} = [f_1, f_2, \alpha, \beta, \gamma, d_u, d_v, a_r].
\end{equation}
The browser also supports a Brown-Conrady model, called \texttt{optmod 20} in
the interface, with radial and tangential coefficients:
\begin{equation}
  \mathbf{p}_{BC} =
  [f_1, f_2, \alpha, \beta, \gamma, d_u, d_v,
   k_1, k_2, k_3, p_1, p_2].
\end{equation}

\begin{table}[h]
\centering
\begin{tabular}{lll}
\toprule
Model & Use case & Notes \\
\midrule
AIDA optmod 2 & all-sky and fisheye cameras & \(\sin(a_r\theta)\) radial law \\
AIDA optmod 12 & flexible all-sky radial law & tangent/sine branch by sign \\
Brown-Conrady & phone and ordinary lenses & radial plus tangential distortion \\
\bottomrule
\end{tabular}
\caption{Common model choices in the browser calibrator.}
\end{table}

\section{Star Detection}

Automatic detection starts from a grayscale representation of the input image.
The detector estimates local and global brightness structure, thresholds
candidate peaks, suppresses nearby duplicates, and rejects obviously elongated
or crowded detections. Detected stars are ranked by contrast and flux-like
scores. Painted bad regions are applied as exclusion masks before automatic
pairing, allowing the user to remove spurious detections from buildings,
aurora, hot pixels, or image artifacts.

Manual picking uses a different path. The user clicks near a star and the tool
extracts a local image patch. A smoothed density estimate is computed on an
interpolated grid, and the local peak is used as the subpixel image coordinate.
This gives the user a controlled way to create high-confidence calibration
pairs even when automatic detection is confused.

\section{Asterism Matching}

The browser contains two matching strategies. The projected matcher uses the
current lens model to project catalog stars into the image and then performs
nearest-neighbor matching with robust filtering. This works well after the
model is already roughly aligned.

The blind spherical matcher is used for bootstrapping. Image detections are
pre-flattened into approximate unit vectors using a grid of candidate focal
scales and radial parameters. Catalog stars visible at the entered time and
site are converted to local sky vectors. Triangular asterisms are represented by
scale-invariant side-ratio signatures and stored in a two-dimensional KD-tree.
For each image triangle, nearby catalog triangle signatures produce candidate
rotations. Candidate rotations are scored by how many additional stars agree
within an angular threshold. The best rotation seeds the lens model and is then
refined by nonlinear fitting.

The current implementation is deliberately bounded for interactivity. It does
not attempt to reproduce the full astrometry.net global quad-index architecture
inside the browser. Instead it uses a Yale bright-star catalog, local visible
stars, and capped triangle sets appropriate for wide-field cameras.

\section{Fitting}

Once a set of star pairs exists, the objective function is the sum of squared
pixel residuals:
\begin{equation}
  J(\mathbf{p}) = \sum_i
  \left[
    \left(x_i - \hat{x}_i(\mathbf{p})\right)^2
    +
    \left(y_i - \hat{y}_i(\mathbf{p})\right)^2
  \right].
\end{equation}
Here \((x_i,y_i)\) is the measured image position and
\((\hat{x}_i,\hat{y}_i)\) is the catalog star projected through the selected
camera model.

Two optimizers are available. The robust Nelder-Mead path tries many jittered
starts, trims bad residuals, and is useful when the initial model is uncertain
or a pair is wrong. The Levenberg-Marquardt path uses numerical derivatives and
is faster near a good solution. The lucky workflow combines these in sweeps:
bright reliable stars are fit first, then additional automatic stars are added
only if they do not substantially increase the root-mean-square error.

\section{Residual Inspection and Undo}

Residual view draws the displacement between each picked image coordinate and
the fitted catalog projection. The application highlights a suggested outlier
based on deviation from the median residual pattern. Accepted fits and
automatic pairing batches are stored on an undo stack, so the user can safely
try automatic expansion and return to the previous state.

\section{Test Cases and Reproducibility}

When the app is served by the local Node server, a calibrated image can be
submitted as a test case. A test case stores image metadata, timestamp,
location, optical model, fitted parameters, matched stars, residuals, and
painted bad star-finder positions. JavaScript unit tests cover coordinate
conversion, camera model behavior, star matching utilities, export generators,
and selected synthetic asterism cases. Heavier image-based tests are skipped by
default and are intended for explicit full-test runs.

\section{Exports}

The browser can export the fitted parameter vector and mapper code in Python,
Julia, C, or MATLAB. The exported code implements the same forward mapping used
by the browser, which maps azimuth and elevation to image pixels for the
selected optical model. This makes the browser calibration directly usable in
downstream AIDA analysis scripts.

\section{Limitations}

The automatic workflow is best suited to wide fields of view above about
30 degrees and scenes containing multiple stars brighter than magnitude 7.
Aurora, buildings, clouds, saturation, and non-star point sources can still
confuse the star detector and asterism matcher. The browser implementation also
uses bounded local data structures for responsiveness; for arbitrary narrow
field blind plate solving, a full astrometry.net-style quad index remains a
better architecture.

\section{Conclusion}

The AIDA browser calibrator provides a practical, inspectable workflow for
turning wide-field star images into calibrated camera models. It combines
manual control with automatic assistance, keeps the calibration state visible,
and exports model code that can be used outside the browser. Version 0.2.0 is a
checkpoint release that stabilizes the lucky fitting workflow, bad detection
masking, undo behavior, and documentation needed for further development.

\section*{Acknowledgments}

This software and article draft were prepared by Juha Vierinen, Bjorn
Gustavsson, and Codex. The implementation builds on the AIDA optical model
conventions and practical calibration workflows developed for all-sky camera
analysis.

\end{document}
